% ---------------------------------------------------------------
\section{Limitations and Future Work}\label{sec:limitations}
% ---------------------------------------------------------------

The Bailout Protocol's guarantees --- SAFETY, LIVENESS, and ACCOUNTABILITY --- are derived from structural invariants enforced by the \fact{} layer's pre-commit hooks. These invariants hold unconditionally within the system's stated assumptions. This section identifies the environmental constraints that bound the protocol's operational scope, classifies each as fundamental or addressable, and charts the research agenda motivated by each. Where a limitation is fundamental, we explain why it cannot be resolved within the current architectural frame. Where it is addressable, we outline the resolution path under Future Work.

% --- 6.1 Limitations ---
\subsection{Limitations}\label{sec:limitations:items}

\subsubsection{Human Response Latency}
\label{sec:lim:latency}

The LIVENESS guarantee is bounded above by the human-response timeout $T_{\text{human}}$, defined in the state-machine specification (Section~\ref{sec:bailout-protocol}). When the designated human anchor $h(n)$ fails to respond within $T_{\text{human}}$, the state machine transitions to \texttt{TIMEOUT} and the originating node enters a quiescent error state. SAFETY is preserved --- consequential action stalls --- but the exception remains unresolved and the system remains in limbo until $h(n)$ regains availability.

The constraint is compounded in globally distributed deployments where $h(n)$ occupies a different timezone, regulatory jurisdiction, or operational context. The protocol currently provides no mechanism to distinguish between an unavailable and an unresponsive anchor, a distinction with materially different operational implications. An unavailable anchor (temporary loss of access, medical emergency, travel) may return shortly and resolve the escalation normally; an unresponsive anchor (cognitive overload, disengagement, or adversarial compromise of the anchor's authentication credentials) represents a categorically different failure that may require pool reassignment or security intervention. Conflating these two states delays appropriate incident response.

Time-sensitive domains --- clinical decision support, real-time infrastructure control, autonomous transport, financial transaction settlement --- cannot guarantee latency-bounded escalation without modifications to the anchor designation model. In these domains, $T_{\text{human}}$ is often smaller than the minimum time required for a human to formulate a calibrated directive under novel exception conditions, creating a structural tension between the timeliness and the quality of the human response. This tension is not resolvable by tuning $T_{\text{human}}$ alone; it requires either domain-specific exemption criteria (which weaken the accountability guarantee) or the adaptive timeout and anchor pool mechanisms described in Section~\ref{sec:future-work}.

We classify the latency limitation as \emph{fundamental under the single-anchor model} and \emph{addressable under the distributed anchor pool model}. The fixed $T_{\text{human}}$ ceiling is a hard bound of the current specification; adaptive variants are a research direction, not a solved mechanism.

\subsubsection{Adversarial Conditions}
\label{sec:lim:adversarial}

The Compulsory Human Escalation theorem (Section~\ref{sec:guarantees-theorem}) assumes a computationally intact parent-pointer chain $\alpha(n)$ at the moment of escalation and a non-adversarial communication substrate. Both assumptions are violated under active adversarial conditions. A compromised parent node $p$ in the propagation chain may exhibit Byzantine behavior: it may acknowledge escalation messages without forwarding them, selectively drop \texttt{BAIL\_PENDING} notifications, or inject fabricated acknowledgments that satisfy the \fact{} layer's entry format while masking the absence of a genuine human response. The parent pointer $\alpha(n)$ remains immutable (the \fact{} layer enforces this), but the node at that pointer position can behave arbitrarily.

The Pathway Health Registry (Section~\ref{sec:bailout-state-machine}) partially mitigates degraded parent nodes by excluding them from pathway selection after a configurable health-check threshold. However, the health check is a network operation --- a probe message and its acknowledgment --- that a sufficiently capable adversary controlling the communication substrate can interfere with or fabricate. The adversary does not need to compromise any individual node; controlling the network transport layer suffices to make the health registry report false positives.

The protocol provides no cryptographic attestation of propagation-path integrity: it cannot guarantee that \texttt{escalation-sent} Forensic Ledger entries correspond to genuinely received and forwarded messages rather than locally fabricated acknowledgments. This gap is fundamental when the adversary controls the network layer rather than individual nodes. Closing it requires a cryptographic propagation proof --- a chain of cryptographic receipts from each node on the path, each attesting to receipt from the previous node --- which the current specification does not provide. Generating such a proof adds per-hop latency and storage overhead proportional to path length, introducing a tradeoff between attestation strength and escalation timeliness.

The diagnostic context attached to each escalation --- exception classification, Bounds Engine reasoning trace, and resolution history of prior occurrences --- additionally exposes an information surface that can be exploited. A compromised node in the propagation chain can observe this context before the escalation reaches its intended terminus. Exception classifications reveal the Boundaries Layer's detection logic; resolution histories expose which failure patterns have and have not been seen before; the Bounds Engine's reasoning trace may expose sensitive business logic, spawn-tree structure, or the principal's operational patterns. An adversary who systematically probes the system to trigger specific exception classes can use the diagnostic responses to build a detailed map of the system's internal architecture.

We classify the adversarial-condition gap as \emph{fundamental under the current single-path propagation model} and \emph{addressable under a multi-path propagation model with cryptographic attestation}. Formal verification of the protocol under Byzantine fault conditions is a prerequisite for any deployment in adversarial environments (Section~\ref{sec:fw:verification}).

\subsubsection{Scalability}
\label{sec:lim:scalability}

The protocol's accountability model is structurally single-threaded at the human anchor layer: each node $n$ designates a single anchor $h(n)$, and all unresolved exceptions route to that one principal. When $N$ simultaneous nodes each trigger the Bailout Protocol, $h(n)$ receives $N$ independent escalations with full diagnostic context. The protocol provides no mechanism for prioritization, batching, automatic triage, or selective compression of concurrent exceptions. All escalations compete equally for the anchor's attention.

The problem is not merely throughput. The single-anchor model creates a cognitive overload condition at high concurrency: even if $h(n)$ could theoretically process $N$ escalations simultaneously, human cognitive bandwidth for exception evaluation is bounded. The protocol provides no schema for comparative risk assessment across concurrent escalations: a low-severity capability-bound violation, a medium-severity timeout, and a high-severity norm violation carry identical escalation formats in the current specification. The anchor must perform triage without structural support from the protocol itself. This is a significant operational burden in deployments with large spawn trees or high exception rates.

Additionally, the forensic record grows linearly with the number of exceptions processed across all nodes. The SHA-256 hash-chained Forensic Ledger imposes storage and verification overhead at scale. Hash-chain verification time is $O(M)$ in the number of entries $M$, meaning that verification of a large ledger becomes progressively more expensive. The current specification provides no tiering, summarization, selective pruning, or checkpoint-based short-circuiting strategy for routine low-severity exceptions. Long-lived systems with high exception throughput will eventually accumulate Forensic Ledgers whose verification cost is non-trivial.

We classify scalability as \emph{addressable through architectural extensions}: distributed anchor pools address the triage bottleneck (Section~\ref{sec:fw:anchor-pools}), and tiered Forensic Ledger management is a natural extension of the existing hash-chain structure that can be specified independently.

\subsubsection{Exploitation Risk}
\label{sec:lim:exploitation}

A motivated adversary who can induce exceptions at a target \bxthree{} node can mount a denial-of-service escalation attack: by repeatedly triggering condition TC, the adversary forces repeated \texttt{BOUNDED} states and eventual \texttt{BAIL\_PENDING} transitions at each targeted node. Each transition consumes the human anchor's evaluation bandwidth, generates Forensic Ledger entries, and may force the node into quiescent error state. The protocol provides no rate-limiting mechanism based on the source, pattern, or temporal density of exception triggers. The anchor has no structural way to distinguish between a node experiencing genuine instability and a node under targeted exception-induction attack.

This risk is particularly acute because TC depends in part on the Bounds Engine's self-assessed confidence threshold $\tau$. An adversary who crafts inputs that systematically reduce the Bounds Engine's confidence below $\tau$ can induce pathological escalation from a well-designed node, even one whose reasoning systems are fundamentally sound. The Bounds Engine's trigger evaluation is treated as a black box: the protocol consumes the resulting trigger signal but provides no external auditability of TC firing rates per node, per input class, or per unit time. A node that fires TC at a high rate generates escalating ledger entries and consumes anchor attention without producing any actionable signal about genuine instability.

While the \fact{} layer prevents machine actors from suppressing escalations once generated, it does not constrain the \emph{generation} of escalation conditions. This asymmetry is deliberate --- constraining generation would require the protocol to make judgmental distinctions about the source and validity of inputs, which would itself introduce new exploitation surfaces --- but it leaves the denial-of-service attack vector unmitigated by design. The correct mitigation is rate-limiting at the anchor interface layer (beyond the protocol's current scope) and canonical anchor-pool triage (which distributes the per-anchor denial-of-service cost across multiple principals).

We classify the exploitation risk as \emph{addressable through rate-limiting extensions at the anchor interface} and \emph{partially mitigated by the distributed anchor pool architecture}, but not fully resolved by the protocol alone.

% --- 6.2 Future Work ---
\subsection{Future Work}\label{sec:future-work}

The limitations above define a structured research agenda organized around three axes: formal verification of the protocol's guarantees under adversarial conditions, adaptive adjustment of timing parameters to operational conditions, and distribution of the human anchor function to eliminate single points of failure. Each axis is independent in its technical scope but interdependent in its practical deployment implications.

\subsubsection{Formal Verification}
\label{sec:fw:verification}

The Bailout Protocol's invariants (INV-1 through INV-3) and the Compulsory Human Escalation theorem are currently stated as pen-and-paper propositions derived from the state-machine definition and \fact{} layer enforcement semantics. While the reasoning is internally consistent, pen-and-paper proofs are subject to subtle logical gaps that mechanized verification is designed to eliminate. A natural next step is mechanized formal verification using TLA+\ \cite{lamport2019tlaplus} or Coq \cite{coq2024}.

Target verification properties include: (i) the \textit{Compulsory Human Escalation theorem}, expressed as a temporal logic formula asserting that for all nodes $n$ and all exception states $e$, the protocol reaches $h(n)$ within a bounded number of transitions; (ii) INV-1, asserting that no node can be in an active-bailout state without the state machine being in \texttt{ESCALATING}; (iii) INV-2, asserting the timeout-forced nature of the \texttt{BAIL\_PENDING} $\rightarrow$ \texttt{ESCALATING} transition; and (iv) INV-3, asserting that the \fact{} layer is the sole source of truth for state machine transition authorization. Mechanizing INV-3 is particularly important: it establishes the separation between the protocol's enforcement (structural) and the agent's reasoning (arbitrary), which is the central correctness claim of the architecture.

Mechanized verification would also enable rigorous analysis of the protocol under Byzantine fault conditions. A formal model could establish the conditions under which the Compulsory Human Escalation theorem continues to hold when up to $f$ parent nodes in the propagation chain exhibit arbitrary failure modes. This requires extending the spawn-tree model to include redundant propagation pathways and majority-voting semantics at the \fact{} layer --- the same extension needed for distributed human anchor pools (Section~\ref{sec:fw:anchor-pools}). Establishing the formal relationship between $f$, the number of redundant paths, and the继续保持 of the theorem's guarantees would provide a principled basis for configuring pool architectures in adversarial deployments.

A TLA+ specification of the \bxthree{} spawn tree and Bailout Protocol state machine is planned as a companion open-reference artifact for researchers and practitioners. This artifact would serve both as a verification substrate and as a reference implementation for practitioners implementing the protocol in non-\bxthree{} architectures.

\subsubsection{Adaptive Timeout}
\label{sec:fw:adaptive-timeout}

The fixed timeout parameters $T_{\text{bounds}}$, $T_{\text{prop}}$, $T_{\text{cap}}$, and $T_{\text{human}}$ are provisioned at node initialization and held constant for the node's operational lifetime. This design simplifies implementation and makes timing behavior statically predictable, which is advantageous for certification and audit contexts. However, static timeouts cannot respond to operational realities: network latency varies with time of day, geographic dispersion introduces persistent baseline delays, compute load variations affect processing time inside node reasoning cycles, and anchor availability varies with timezone and work schedule.

We intend to investigate an adaptive timeout extension in which each timeout parameter is replaced by a feedback-controlled function that adjusts based on observed historical values. Concretely, let $T_{\text{prop}}^{\text{obs}}(n)$ denote the empirically observed propagation latency from node $n$ to $h(n)$ across the last $k$ successful escalations. An adaptive $T_{\text{prop}}$ could be set as a configured percentile (e.g., 95th) of $T_{\text{prop}}^{\text{obs}}$ plus a safety margin, recomputed on a rolling-window basis. This would accommodate diurnal network latency variation, geographic dispersion of nodes, and timezone effects on anchor availability without manual re-provisioning.

The primary risk of adaptive timeouts is adversarial manipulation of observed latencies to compress the effective timeout below the level needed for a human to formulate a meaningful response. An adversary controlling the network substrate could systematically reduce observed latency reports, causing the adaptive function to converge on dangerously low timeout values. This risk is mitigated by establishing hard lower bounds on each adaptive parameter: $T_{\text{human}}$ must always exceed a provisioned floor $T_{\text{human}}^{\min}$ regardless of observed values, and the \fact{} layer enforces this floor as an invariant. Formal analysis of feedback safety under adversarial manipulation --- specifically, proving that no adversary can cause the adaptive function to converge below its floor by manipulating observable latencies --- is a required deliverable of this research track and is interdependent with the formal verification agenda (Section~\ref{sec:fw:verification}).

\subsubsection{Distributed Human Anchor Pools}
\label{sec:fw:anchor-pools}

The single-anchor designation $h(n)$ is the root cause of three of the four limitations described above: the latency ceiling (unavailable single anchor), the scalability bottleneck (single-threaded triage), and the exploitation vulnerability (focused denial-of-service bandwidth). This motivates the investigation of a distributed human anchor pool architecture as the primary structural extension to the protocol.

Under the pool model, rather than designating a single principal $h(n)$ for each node $n$, the protocol designates a pool $H(n) = \{h_1, h_2, \ldots, h_m\}$ of $m$ human anchors, each with authenticated access to the \fact{} layer's escalation endpoint. The protocol requires acknowledgment from any $k$-subset of $H(n)$ to transition from \texttt{ESCALATING} to \texttt{HUMAN\_ACKNOWLEDGED}, providing Byzantine fault tolerance against up to $m - k$ unresponsive or adversarial anchors. The values of $m$ and $k$ are provisioned at node initialization and recorded in the Fact Layer; they are adjustable subject to \fact{} layer immutability constraints (changing the pool requires a new Fact Layer entry, not an overwrite).

The pool architecture addresses all three single-anchor failure modes. Availability improves because the system tolerates the absence of any individual anchor without blocking escalation. Scalability improves because the triage burden is distributed: a high-severity exception can be routed to the first available anchor rather than waiting for a designated principal to become reachable, and multiple anchors can process concurrent escalations in parallel. The denial-of-service attack vector is widened but its per-anchor cost is reduced: an adversary must now exhaust the attention of $k$ independent anchors to achieve the same effect that exhausting a single anchor achieved under the original model.

The pool architecture introduces open design questions that require systematic resolution. First, directive-issuance semantics must handle the case where different anchors in the pool issue conflicting directives for the same escalation. A majority-voting scheme could resolve this --- the \fact{} layer records all submitted directives and the final directive is determined by plurality --- but this requires the \fact{} layer to support multi-author directive entries and introduces non-determinism in the resolution path that must be accounted for in the forensic record. Second, dynamic pool reconfiguration must not invalidate in-flight escalations: an anchor that departs mid-escalation (voluntary exit, credential revocation, adversarial compromise) must not cause the protocol to stall or deadlock. A minimum-active-anchors threshold, below which no new escalations can be initiated but in-flight escalations continue with the remaining anchors, is one structural approach. Third, pool membership at the time of each escalation must be recorded in the Forensic Ledger, and retroactive changes to pool membership must not affect the verifiability of past escalation records --- this is achievable through the Fact Layer's append-only semantics but requires careful specification. Fourth, pool sizing and $k$-threshold calibration for different deployment contexts (high-stakes vs. routine operations) is an operational question that requires empirical validation.

% ---------------------------------------------------------------
\section{Conclusion}\label{sec:conclusion}
% ---------------------------------------------------------------

The Bailout Protocol addresses a foundational problem in accountable autonomous systems: the absence of an architecturally enforced mechanism that compels unresolved exception states to reach a designated human principal, bypassing all machine actors in the intervening hierarchy. We have shown that this compulsion is not achieved through policy or convention --- the two approaches that characterize the majority of HITL systems in current deployment --- but through the structural integration of the protocol's state machine into the \fact{} layer's pre-commit enforcement substrate. The integration makes it impossible for any machine actor, including the agent whose behavior triggered the exception, to suppress, redirect, or substitute the human escalation path. This is the protocol's core claim, and it is enforced architecturally rather than argued operationally.

The protocol makes three concrete contributions. First, it provides a precise operational semantics for exception escalation in bounded-autonomy agentic systems, expressed as a deterministic state machine $\mathcal{B}$ with formally specified invariants (INV-1 through INV-3) that encode the protocol's core safety guarantees and a clearly defined trigger condition TC that governs activation. The state-machine specification is self-contained: its properties are derived from its own structure and from the \fact{} layer's enforcement semantics, not from assumptions about agent behavior. Second, it establishes the \textit{Compulsory Human Escalation theorem}: for any node $n$ in a \bxthree{} spawn tree and any exception state $e$ at $n$, the protocol's propagation path $C(n)$ reaches $h(n)$ in finite time, independent of the reasoning strategy employed by any machine actor in the chain. The theorem's guarantee is structural, not probabilistic; it follows from the immutability of the parent pointer and the finiteness of the spawn tree depth, both enforced by the \fact{} layer. Third, it demonstrates that the upstream accountability guarantee of the \bxthree{} Framework --- that systems fail \emph{upward} into human consciousness, never \emph{downward} into autonomous chaos --- is not a design aspiration but an architecturally enforced operational property, made compulsory by the three-layer model and the pre-commit enforcement mechanism.

The \bxthree{} Framework's role is constitutive of all three contributions, not merely incidental. The three-layer model provides the structural substrate on which the protocol operates: the \purpose{} layer defines the human anchor $h(n)$ and the node's authorized operational scope; the \bounds{} layer houses the Bounds Engine that detects and classifies constraint violations and evaluates TC; and the \fact{} layer enforces the state machine transitions that drive the protocol, records every escalation event in the immutable Forensic Ledger, and prevents any machine actor from intercepting or modifying the escalation path. Without this three-layer separation, the Bailout Protocol would reduce to a discretionary interface whose enforcement depended on the correct behavior of the very machine actors it is designed to govern. The \bxthree{} Framework is not a container for the Bailout Protocol; it is the mechanism by which the protocol becomes compulsory.

We acknowledge that the protocol operates within real constraints: human response latency imposes a practical upper bound on escalation timeliness that is fundamental under the single-anchor model; adversarial network conditions can compromise the propagation chain in ways that the current single-path specification cannot detect or resist; single-anchor designation limits scalability and creates a focused exploitation surface; and the escalation interface, by design, exposes diagnostic context that a compromised node in the propagation chain can observe. These are not reasons to abandon the approach --- they are the precise targets of the formal verification, adaptive timeout, and distributed human anchor pool research programs outlined in Section~\ref{sec:future-work}. Each limitation corresponds to an open research question with a clear resolution path, and each resolution path is architecturally compatible with the protocol's existing invariants.

The broader significance of the Bailout Protocol extends beyond the \bxthree{} Framework's specific architecture. The protocol demonstrates that the upstream accountability guarantee --- a long-standing design principle in human-in-the-loop AI literature, traceable to Wiener's original conception of the human sebagai last resort \cite{wiener1954human} --- can be rendered formally compulsory through architectural enforcement rather than left to operational convention. The structural principles underlying the protocol --- state-machine embedding in a tamper-evident enforcement layer, immutable human-anchor designation, and deterministic timeout-forced escalation --- are applicable to accountable autonomous system design more broadly, independent of the specific reasoning architecture or domain of deployment. Any autonomous system that issues consequential actions without requiring human review of the decision process is, by design, incapable of failing safely. The Bailout Protocol shows what failing safely actually requires: not a human in the loop, but a structural compulsion that makes the human-in-the-loop unconditional.
